\documentclass[a4paper,11pt,reqno]{amsart}
\usepackage[margin=1.2in,marginparwidth=1.5cm,marginparsep=0.5cm]{geometry}
\usepackage{graphicx}
\usepackage{amsmath,amssymb,amsthm}
\usepackage{mathtools}

\title{Solucionário do Simulado de ``Introdução às Equações Diferenciais''  IMPA-TECH 2026-01}

\begin{document}

\maketitle

\section*{Solucionário}

\begin{enumerate}

\item Resolva os seguintes problemas de valor inicial. Em cada item, indique explicitamente o maior intervalo em que a solução está definida.

\begin{enumerate}

\item
\[
    y'+2y=e^{-t},
    \qquad
    y(0)=1.
\]

Esta é uma equação linear de primeira ordem. O fator integrante é
\[
    \mu(t)=e^{\int 2\,dt}=e^{2t}.
\]
Multiplicando a equação por $e^{2t}$, obtemos
\[
    e^{2t}y'+2e^{2t}y=e^t,
\]
isto é,
\[
    \left(e^{2t}y(t)\right)'=e^t.
\]
Integrando,
\[
    e^{2t}y(t)=e^t+C.
\]
Portanto
\[
    y(t)=e^{-t}+Ce^{-2t}.
\]
Usando a condição inicial $y(0)=1$,
\[
    1=1+C,
\]
logo $C=0$. Assim,
\[
    \boxed{y(t)=e^{-t}.}
\]
A solução está definida para todo $t\in\mathbb{R}$. Portanto o maior intervalo de existência é
\[
    \boxed{\mathbb{R}.}
\]

\item
\[
    y'=t(1+y^2),
    \qquad
    y(0)=0.
\]

A equação é separável:
\[
    \frac{dy}{1+y^2}=t\,dt.
\]
Integrando,
\[
    \arctan y=\frac{t^2}{2}+C.
\]
Como $y(0)=0$, temos
\[
    \arctan(0)=0=C.
\]
Logo
\[
    \arctan y=\frac{t^2}{2},
\]
e portanto
\[
    \boxed{y(t)=\tan\left(\frac{t^2}{2}\right).}
\]

Agora precisamos determinar o maior intervalo contendo $0$ no qual essa expressão está definida e resolve o PVI. A função tangente tem assíntotas quando seu argumento é da forma
\[
    \frac{\pi}{2}+k\pi,
    \qquad k\in\mathbb{Z}.
\]
Perto de $t=0$, o primeiro problema ocorre quando
\[
    \frac{t^2}{2}=\frac{\pi}{2},
\]
isto é,
\[
    t^2=\pi.
\]
Portanto o maior intervalo contendo $0$ é
\[
    \boxed{(-\sqrt{\pi},\sqrt{\pi}).}
\]

\item
\[
    y'=\frac{t+y}{t},
    \qquad
    y(1)=0,
    \qquad
    t>0.
\]

Podemos reescrever a equação como
\[
    y'=1+\frac{y}{t}.
\]
A expressão $\frac{y}{t}$ sugere a mudança de variável
\[
    z(t)=\frac{y(t)}{t},
    \qquad\text{isto é,}\qquad
    y(t)=tz(t).
\]
Então
\[
    y'(t)=z(t)+tz'(t).
\]
Substituindo na equação,
\[
    z+tz'=1+z.
\]
Cancelando $z$ dos dois lados,
\[
    tz'=1.
\]
Como $t>0$,
\[
    z'=\frac{1}{t}.
\]
Integrando,
\[
    z(t)=\ln t+C.
\]
Logo
\[
    y(t)=t(\ln t+C).
\]
Usando $y(1)=0$,
\[
    0=1(\ln 1+C)=C.
\]
Portanto
\[
    \boxed{y(t)=t\ln t.}
\]
Como a própria equação só está considerada para $t>0$, e a solução está definida em todo esse intervalo, o maior intervalo de existência contendo $1$ é
\[
    \boxed{(0,+\infty).}
\]

\end{enumerate}


\item Considere o problema de valor inicial
\[
    \frac{dy}{dt}=3y^{2/3},
    \qquad
    y(t_0)=y_0.
\]

\begin{enumerate}

\item No ponto $(t_0,y_0)=(0,1)$, temos o PVI
\[
    y'=3y^{2/3},
    \qquad
    y(0)=1.
\]
A função
\[
    f(t,y)=3y^{2/3}
\]
é contínua em todo o plano. Além disso,
\[
    \frac{\partial f}{\partial y}(t,y)=2y^{-1/3}
\]
é contínua em qualquer região que não contenha $y=0$. Como o ponto inicial é $(0,1)$, existe uma vizinhança desse ponto onde $f$ e $\partial f/\partial y$ são contínuas. Portanto, pelo teorema de existência e unicidade, existe uma única solução local passando por $(0,1)$.

Vamos encontrá-la explicitamente. Para $y\neq 0$, podemos escrever
\[
    \frac{dy}{dt}=3y^{2/3}.
\]
Como
\[
    \frac{d}{dt}\left(y^{1/3}\right)
    =
    \frac{1}{3}y^{-2/3}y',
\]
obtemos
\[
    \frac{d}{dt}\left(y^{1/3}\right)=1.
\]
Logo
\[
    y^{1/3}=t+C.
\]
Portanto
\[
    y(t)=(t+C)^3.
\]
Usando $y(0)=1$,
\[
    1=C^3,
\]
e então $C=1$. Assim, a solução local é
\[
    \boxed{y(t)=(t+1)^3.}
\]

Em particular, em qualquer intervalo suficientemente pequeno contendo $0$, essa solução é única.

Observação: se perguntarmos por extensões globais em intervalos maiores que atravessam o ponto $t=-1$, a não unicidade pode reaparecer depois que a solução atinge $y=0$. Por exemplo,
\[
    y(t)=(t+1)^3
\]
e
\[
    y(t)=
    \begin{cases}
        0, & t\leq -1,\\
        (t+1)^3, & t\geq -1,
    \end{cases}
\]
são ambas soluções globais que satisfazem $y(0)=1$. Isso não contradiz a unicidade local perto de $t=0$.

\item No ponto $(t_0,y_0)=(0,0)$, temos
\[
    y'=3y^{2/3},
    \qquad
    y(0)=0.
\]
A solução identicamente nula,
\[
    \boxed{y(t)\equiv 0,}
\]
é uma solução.

Por outro lado,
\[
    \boxed{y(t)=t^3}
\]
também é solução, pois
\[
    y'(t)=3t^2,
\]
enquanto
\[
    3y(t)^{2/3}=3(t^3)^{2/3}=3t^2.
\]
Além disso,
\[
    y(0)=0.
\]
Portanto o PVI não tem solução única.

Na verdade, há infinitas soluções. Para cada $a\geq 0$, a função
\[
    y_a(t)=
    \begin{cases}
        0, & t\leq a,\\
        (t-a)^3, & t\geq a,
    \end{cases}
\]
é solução e satisfaz $y_a(0)=0$.

\item A diferença entre os dois casos está na hipótese de unicidade do teorema de existência e unicidade.

A função
\[
    f(t,y)=3y^{2/3}
\]
é contínua em todo lugar, portanto o teorema de existência garante existência local para qualquer condição inicial.

No entanto, a derivada parcial em relação a $y$ é
\[
    \frac{\partial f}{\partial y}(t,y)=2y^{-1/3},
\]
que não é contínua em $y=0$.

No ponto $(0,1)$, estamos longe de $y=0$, então $f$ é localmente Lipschitz em $y$ e a solução local é única.

No ponto $(0,0)$, a hipótese de unicidade falha. Isso não prova automaticamente que a solução não é única, mas abre a possibilidade de não unicidade. De fato, como vimos no item anterior, há várias soluções passando por $(0,0)$.

\end{enumerate}


\item Considere a equação diferencial
\[
    t^2y''-2ty'+2y=t^3,
    \qquad
    t>0.
\]

\begin{enumerate}

\item A equação homogênea associada é
\[
    t^2y''-2ty'+2y=0.
\]
Para $y_1(t)=t$, temos
\[
    y_1'(t)=1,
    \qquad
    y_1''(t)=0.
\]
Substituindo,
\[
    t^2\cdot 0-2t\cdot 1+2t=0.
\]
Logo
\[
    \boxed{y_1(t)=t}
\]
é solução da equação homogênea.

\item Procuramos soluções da forma
\[
    y(t)=t^m.
\]
Então
\[
    y'(t)=mt^{m-1},
    \qquad
    y''(t)=m(m-1)t^{m-2}.
\]
Substituindo na equação homogênea,
\[
    t^2m(m-1)t^{m-2}
    -
    2tmt^{m-1}
    +
    2t^m
    =
    0.
\]
Isto é,
\[
    \left[m(m-1)-2m+2\right]t^m=0.
\]
Logo o polinômio característico é
\[
    m(m-1)-2m+2=m^2-3m+2=(m-1)(m-2).
\]
Assim,
\[
    m=1
    \qquad\text{ou}\qquad
    m=2.
\]
Já conhecíamos a solução $t$. A segunda solução é
\[
    \boxed{y_2(t)=t^2.}
\]
Como $t$ e $t^2$ não são múltiplas uma da outra em $(0,+\infty)$, elas são linearmente independentes.

\item A solução geral da homogênea é
\[
    y_h(t)=C_1t+C_2t^2.
\]

Para encontrar uma solução particular da equação não homogênea
\[
    t^2y''-2ty'+2y=t^3,
\]
tentamos
\[
    y_p(t)=At^3.
\]
Então
\[
    y_p'(t)=3At^2,
    \qquad
    y_p''(t)=6At.
\]
Substituindo,
\[
    t^2(6At)-2t(3At^2)+2At^3
    =
    6At^3-6At^3+2At^3
    =
    2At^3.
\]
Queremos que isso seja igual a $t^3$, logo
\[
    2A=1,
    \qquad
    A=\frac12.
\]
Assim,
\[
    y_p(t)=\frac12t^3.
\]
Portanto a solução geral é
\[
    \boxed{y(t)=C_1t+C_2t^2+\frac12t^3.}
\]

\item Devemos impor
\[
    y(1)=0,
    \qquad
    y'(1)=1.
\]
Temos
\[
    y(t)=C_1t+C_2t^2+\frac12t^3
\]
e
\[
    y'(t)=C_1+2C_2t+\frac32t^2.
\]
Logo
\[
    y(1)=C_1+C_2+\frac12=0,
\]
isto é,
\[
    C_1+C_2=-\frac12.
\]
Além disso,
\[
    y'(1)=C_1+2C_2+\frac32=1,
\]
isto é,
\[
    C_1+2C_2=-\frac12.
\]
Subtraindo a primeira equação da segunda, obtemos
\[
    C_2=0.
\]
Então
\[
    C_1=-\frac12.
\]
Portanto
\[
    \boxed{
    y(t)=\frac12t^3-\frac12t.
    }
\]

Como a equação está definida para $t>0$, e o ponto inicial é $t=1$, o maior intervalo natural de existência é
\[
    \boxed{(0,+\infty).}
\]

\end{enumerate}


\item Seja $a:[0,1]\to(0,+\infty)$ uma função contínua, e seja $\lambda\in\mathbb{R}$. Considere a equação
\[
    x''+\lambda a(t)x=0.
\]
Suponha que exista uma solução não identicamente nula $x$ tal que
\[
    x(0)=0,
    \qquad
    x'(1)=0.
\]

\begin{enumerate}

\item Suponha, por absurdo, que $\lambda\leq 0$ e que exista uma solução não identicamente nula satisfazendo as condições de contorno.

Multiplicamos a equação por $x(t)$:
\[
    x''(t)x(t)+\lambda a(t)x(t)^2=0.
\]
Integramos de $0$ a $1$:
\[
    \int_0^1 x''(t)x(t)\,dt
    +
    \lambda\int_0^1 a(t)x(t)^2\,dt
    =
    0.
\]
Pela fórmula de integração por partes,
\[
    \int_0^1 x''(t)x(t)\,dt
    =
    \left[x'(t)x(t)\right]_0^1
    -
    \int_0^1 (x'(t))^2\,dt.
\]
Agora,
\[
    \left[x'(t)x(t)\right]_0^1
    =
    x'(1)x(1)-x'(0)x(0).
\]
Como $x'(1)=0$ e $x(0)=0$, esse termo é nulo. Portanto
\[
    \int_0^1 x''(t)x(t)\,dt
    =
    -
    \int_0^1 (x'(t))^2\,dt.
\]
Substituindo na identidade integral,
\[
    -
    \int_0^1 (x'(t))^2\,dt
    +
    \lambda\int_0^1 a(t)x(t)^2\,dt
    =
    0.
\]
Assim,
\[
    \int_0^1 (x'(t))^2\,dt
    =
    \lambda\int_0^1 a(t)x(t)^2\,dt.
\]
Se $\lambda\leq 0$, o lado direito é menor ou igual a zero, pois
\[
    a(t)>0
    \qquad\text{e}\qquad
    x(t)^2\geq 0.
\]
Por outro lado, o lado esquerdo é sempre maior ou igual a zero. Logo ambos os lados devem ser iguais a zero. Em particular,
\[
    \int_0^1 (x'(t))^2\,dt=0.
\]
Como $(x')^2$ é contínua e não negativa, segue que
\[
    x'(t)=0
    \qquad\text{para todo }t\in[0,1].
\]
Logo $x$ é constante. Como $x(0)=0$, concluímos que
\[
    x(t)\equiv 0,
\]
contradizendo a hipótese de que $x$ não é identicamente nula.

Portanto, se $\lambda\leq 0$, tal solução não pode existir.

\item Do item anterior, concluímos imediatamente que a existência de uma solução não identicamente nula satisfazendo
\[
    x(0)=0,
    \qquad
    x'(1)=0
\]
implica
\[
    \boxed{\lambda>0.}
\]

\item Tomemos
\[
    a(t)\equiv 1
\]
e
\[
    \lambda=\frac{\pi^2}{4}.
\]
A equação fica
\[
    x''+\frac{\pi^2}{4}x=0.
\]
Considere
\[
    x(t)=\sin\left(\frac{\pi t}{2}\right).
\]
Então
\[
    x(0)=\sin 0=0.
\]
Além disso,
\[
    x'(t)=\frac{\pi}{2}\cos\left(\frac{\pi t}{2}\right),
\]
e portanto
\[
    x'(1)=\frac{\pi}{2}\cos\left(\frac{\pi}{2}\right)=0.
\]
A função $x$ não é identicamente nula. Logo este é um exemplo com $\lambda>0$.

\end{enumerate}


\item Seja $I\subset\mathbb{R}$ um intervalo aberto e seja
\[
    L:C^2(I)\to C^0(I)
\]
um operador da forma
\[
    L[y](t)=y''(t)+p(t)y'(t)+q(t)y(t),
\]
onde $p,q:I\to\mathbb{R}$ são funções contínuas.

\begin{enumerate}

\item Queremos mostrar que não existe uma equação linear homogênea de segunda ordem com coeficientes contínuos em todo $\mathbb{R}$,
\[
    y''+p(t)y'+q(t)y=0,
\]
que possua simultaneamente
\[
    y_1(t)=\sin(t^2),
    \qquad
    y_2(t)=\cos(t^2)
\]
como soluções em todo $\mathbb{R}$.

Calculamos primeiro o Wronskiano:
\[
    W(y_1,y_2)(t)
    =
    \begin{vmatrix}
        y_1(t) & y_2(t)\\
        y_1'(t) & y_2'(t)
    \end{vmatrix}
    =
    y_1(t)y_2'(t)-y_1'(t)y_2(t).
\]
Temos
\[
    y_1'(t)=2t\cos(t^2),
    \qquad
    y_2'(t)=-2t\sin(t^2).
\]
Logo
\[
    W(y_1,y_2)(t)
    =
    \sin(t^2)\bigl(-2t\sin(t^2)\bigr)
    -
    \bigl(2t\cos(t^2)\bigr)\cos(t^2).
\]
Portanto
\[
    W(y_1,y_2)(t)
    =
    -2t\sin^2(t^2)-2t\cos^2(t^2)
    =
    -2t.
\]
Assim,
\[
    \boxed{W(y_1,y_2)(t)=-2t.}
\]

Em particular,
\[
    W(y_1,y_2)(0)=0,
\]
mas
\[
    W(y_1,y_2)(t)\neq 0
    \qquad\text{se }t\neq 0.
\]

Suponha, por absurdo, que existam funções contínuas $p,q:\mathbb{R}\to\mathbb{R}$ tais que $y_1$ e $y_2$ sejam soluções da mesma equação
\[
    y''+p(t)y'+q(t)y=0
\]
em todo $\mathbb{R}$.

Pela identidade de Abel, o Wronskiano de duas soluções deve satisfazer
\[
    W(t)=W(t_0)\exp\left(-\int_{t_0}^{t}p(s)\,ds\right).
\]
Tomando $t_0=0$, como $W(0)=0$, obtemos
\[
    W(t)=0\cdot\exp\left(-\int_0^{t}p(s)\,ds\right)=0
\]
para todo $t\in\mathbb{R}$. Portanto o Wronskiano deveria ser identicamente nulo em $\mathbb{R}$.

Mas calculamos que
\[
    W(t)=-2t,
\]
que não é identicamente nulo. Contradição.

Logo não existe tal equação em todo $\mathbb{R}$.

\item O resultado anterior não vale se substituirmos $\mathbb{R}$ por $(0,+\infty)$. De fato, em $(0,+\infty)$ o Wronskiano é
\[
    W(t)=-2t\neq 0.
\]
Assim, as funções $y_1$ e $y_2$ podem formar uma base de soluções de uma equação linear homogênea de segunda ordem com coeficientes contínuos em $(0,+\infty)$.

Vamos encontrar essa equação.

Queremos funções $p,q$ tais que
\[
    y''+p(t)y'+q(t)y=0
\]
tenha
\[
    y_1(t)=\sin(t^2),
    \qquad
    y_2(t)=\cos(t^2)
\]
como soluções.

Pela identidade de Abel,
\[
    W'=-pW.
\]
Como
\[
    W(t)=-2t,
    \qquad
    W'(t)=-2,
\]
temos
\[
    -2=-p(t)(-2t).
\]
Logo
\[
    -2=2tp(t),
\]
isto é,
\[
    \boxed{p(t)=-\frac{1}{t}.}
\]

Agora usamos uma das soluções para determinar $q$. Para
\[
    y_1(t)=\sin(t^2),
\]
temos
\[
    y_1'(t)=2t\cos(t^2)
\]
e
\[
    y_1''(t)=2\cos(t^2)-4t^2\sin(t^2).
\]
Substituindo na equação,
\[
    y_1''+p(t)y_1'+q(t)y_1=0.
\]
Como $p(t)=-1/t$,
\[
    \left(2\cos(t^2)-4t^2\sin(t^2)\right)
    -
    \frac{1}{t}\left(2t\cos(t^2)\right)
    +
    q(t)\sin(t^2)
    =
    0.
\]
Os termos com cosseno cancelam:
\[
    -4t^2\sin(t^2)+q(t)\sin(t^2)=0.
\]
Logo
\[
    q(t)=4t^2.
\]
Portanto a equação é
\[
    \boxed{
        y''-\frac{1}{t}y'+4t^2y=0,
        \qquad
        t>0.
    }
\]

Os coeficientes
\[
    p(t)=-\frac1t,
    \qquad
    q(t)=4t^2
\]
são contínuos em $(0,+\infty)$. Portanto, nesse intervalo, existe uma equação linear homogênea de segunda ordem com coeficientes contínuos que possui $\sin(t^2)$ e $\cos(t^2)$ como soluções.

A mesma fórmula também define uma equação em $(-\infty,0)$, mas não em todo $\mathbb{R}$, pois o coeficiente $-1/t$ não é contínuo em $t=0$.

\end{enumerate}

\end{enumerate}

\end{document}